\documentclass[11pt,letterpaper]{article}

% --- packages -----------------------------------------------------------
\usepackage[margin=1in]{geometry}
\usepackage{amsmath, amssymb, amsthm}
\usepackage{bm}
\usepackage{booktabs}
\usepackage{xcolor}
\usepackage{enumitem}
\usepackage[hidelinks]{hyperref}
\usepackage{graphicx}
\usepackage{caption}
\captionsetup{font=small,labelfont=bf}

\title{\textbf{Tumor--Immune Active Matter}\\
       \large Mathematical and physical model summary}
\author{I-Tsen Tsai \\ {\small \texttt{i3tsai@ucsd.edu}}}
\date{Vibe Coding Active Matter \& Biophysics Hackathon \\ UCSD --- May 23, 2026}

\begin{document}
\maketitle

\begin{abstract}
We model the tumor microenvironment as a multi-species active mixture in a 2D
periodic box: slow self-propelled tumor cells that proliferate and secrete
diffusing chemicals, fast self-propelled CD8 T cells that chemotax up an
attractant and down a suppressant and kill on contact, and macrophages
that carry a continuous M1/M2 polarization driven by local cytokine
signals. The cells are embedded in an extracellular matrix (ECM) density
field that is degraded by tumor-secreted MMPs and that gates cell motion
through drag and pore-size cutoffs. Cadherin-mediated cell--cell
adhesion plus a Vicsek alignment torque is layered onto the tumor pairwise
interaction. The full coupled system --- non-conserved cell dynamics,
multi-scale chemical signaling with opposite signs, overdamped Langevin
self-propulsion, mechanical-pressure-gated proliferation, ECM remodeling,
and intercellular adhesion --- generates a phase diagram with three
biologically meaningful regimes (clearance, dormancy, escape) that map
directly onto the clinical hot/excluded/cold tumor classification. This
document collects every equation, parameter value, and numerical scheme
in one place.
\end{abstract}

\section{Setup}

A 2D periodic square box of side $L$ contains
\begin{itemize}[nosep]
  \item $N_T(t)$ tumor cells with positions $\mathbf{r}_i$ and headings
        $\theta_i$, $i=1,\dots,N_T$;
  \item $N_I(t)$ T cells with positions $\mathbf{R}_j$ and headings
        $\phi_j$, $j=1,\dots,N_I$;
  \item two scalar concentration fields $c_a(\mathbf{x},t)$ (attractant)
        and $c_s(\mathbf{x},t)$ (suppressant) defined on a $G\times G$
        grid covering the box.
\end{itemize}
Tumor cells are non-conserved: born by proliferation, removed by killing.
T cells are conserved (no proliferation or natural death on the
simulation timescale).

\section{Tumor-cell dynamics}

Tumor cells obey overdamped Langevin self-propulsion with soft pairwise
repulsion and rotational diffusion:
\begin{align}
\frac{d\mathbf{r}_i}{dt}
  &= v_T\, \mathbf{n}_i
   + \sum_{j\ne i} \mathbf{F}_{\text{rep}}(\mathbf{r}_i - \mathbf{r}_j)
   + \sum_{k}      \mathbf{F}_{\text{rep}}(\mathbf{r}_i - \mathbf{R}_k)
   + \sqrt{2 D_T}\;\bm{\eta}_i(t), \label{eq:tumor_pos}\\[4pt]
\frac{d\theta_i}{dt}
  &= \sqrt{2 D_R^T}\,\xi_i(t), \label{eq:tumor_theta}
\end{align}
with $\mathbf{n}_i = (\cos\theta_i, \sin\theta_i)$ and unit-variance
white noises $\bm{\eta}_i, \xi_i$ independent across particles and time.
The repulsion is harmonic with cutoff at one cell diameter $\sigma$:
\begin{equation}
\mathbf{F}_{\text{rep}}(\mathbf{d})
   = k_{\text{rep}}\,\max(0,\,\sigma - \|\mathbf{d}\|)\,\hat{\mathbf{d}}.
\label{eq:harmonic}
\end{equation}

\paragraph{Proliferation.}
At every step, each alive tumor cell $i$ attempts to divide independently
with probability $p_{\text{div}}$, conditional on a local-density gate
\begin{equation}
n_i \;\equiv\; \big|\{j\ne i : \|\mathbf{r}_i - \mathbf{r}_j\| < r_{\text{nbr}}\}\big|
                \;<\; n_{\max}.
\label{eq:density_gate}
\end{equation}
If the gate passes, a daughter cell is placed at
$\mathbf{r}_{i'} = \mathbf{r}_i + 0.3\,\sigma\,(\cos\psi,\sin\psi)$
with $\psi\sim\mathcal{U}[0,2\pi)$, and a uniformly random heading.

\section{T-cell dynamics}

T cells are coupled to both chemical fields:
\begin{align}
\boxed{\;
\frac{d\mathbf{R}_j}{dt}
  = v_I\, \mathbf{m}_j
  + \chi_a\, \nabla c_a\big|_{\mathbf{R}_j}
  - \alpha\,  \nabla c_s\big|_{\mathbf{R}_j}
  + \sum_{k\ne j} \mathbf{F}_{\text{rep}}(\mathbf{R}_j - \mathbf{R}_k)
  + \sum_i        \mathbf{F}_{\text{rep}}(\mathbf{R}_j - \mathbf{r}_i)
  + \sqrt{2 D_I}\,\bm{\zeta}_j(t)\;
}
\label{eq:tcell_pos}
\end{align}
\begin{equation}
\frac{d\phi_j}{dt}
   = \sqrt{2 D_R^I}\,\mu_j(t).
\label{eq:tcell_phi}
\end{equation}
The crucial physical content is in the two chemotactic terms:
\begin{itemize}[nosep]
  \item $+\chi_a\,\nabla c_a$ pulls T cells \emph{up} the attractant
        gradient --- toward tumor-secreted signals.
  \item $-\alpha\,\nabla c_s$ pushes T cells \emph{down} the suppressant
        gradient --- away from the tumor's immunosuppressive halo.
\end{itemize}
The coupling $\alpha\equiv\chi_s$ is the immunosuppression strength
and is one of the two control parameters of the phase diagram. The other is
$\rho_I\equiv N_I(0)$.

\section{Reaction--diffusion fields}

Both fields obey a diffusion--source--decay PDE with the tumor density
$\rho_T(\mathbf{x},t)$ as the source:
\begin{align}
\frac{\partial c_a}{\partial t}
   &= D_a\,\nabla^2 c_a + s_a\,\rho_T - \lambda_a\, c_a,
   \label{eq:rd_a}\\[4pt]
\frac{\partial c_s}{\partial t}
   &= D_s\,\nabla^2 c_s + s_s\,\rho_T - \lambda_s\, c_s.
   \label{eq:rd_s}
\end{align}
The defining asymmetry is $D_a \gg D_s$: attractant diffuses far, suppressant
stays local. This separates spatial scales --- T cells can sense the tumor
from across the box (long-range $c_a$), but only feel repulsion close in
(short-range $c_s$). It is the competition between these two ranges, with
opposite signs in the T-cell equation~\eqref{eq:tcell_pos}, that generates
the phase diagram.

\paragraph{Steady-state ranges.}
Linearizing about a localized tumor source, the two fields have
characteristic length scales
\begin{equation}
\ell_a = \sqrt{D_a/\lambda_a} \approx 7\sigma, \qquad
\ell_s = \sqrt{D_s/\lambda_s} \approx 2.2\sigma,
\label{eq:lengths}
\end{equation}
i.e.\ a ratio of $\approx 3.2$ at our parameter values.

\paragraph{Cloud-in-cell deposition.}
The tumor density on the grid is computed by bilinear deposition. For each
tumor cell at fractional grid coordinates $(g_x, g_y)$ with integer parts
$(I_x, I_y)$ and fractional parts $(f_x, f_y)$:
\begin{equation}
\rho_T(I_x{+}p, I_y{+}q) \mathrel{+}=
   \frac{w_{p,q}}{(\Delta x)^2}, \qquad
w_{p,q} = (1{-}f_x)^{1-p} f_x^{\,p}\,(1{-}f_y)^{1-q} f_y^{\,q},
\label{eq:cic}
\end{equation}
for $(p,q)\in\{0,1\}^2$. Total mass is preserved exactly.

\section{Killing rule}

For each alive T cell $j$, the rule iterates as follows:
\begin{enumerate}[nosep,leftmargin=*]
  \item Find the nearest alive tumor cell $i^\star(j)$ under minimum-image
        distance.
  \item If $\|\mathbf{R}_j - \mathbf{r}_{i^\star}\| < r_{\text{kill}}$
        \emph{and} $U \sim \mathcal{U}[0,1) < p_{\text{kill}}$, remove
        tumor cell $i^\star$.
\end{enumerate}
The iteration is over T cells (\emph{not} over tumor cells), so the kill
rate scales with $N_I$, not $N_T$. At most one kill per T cell per step.

\section{Numerical methods}

\subsection{Euler--Maruyama (particles)}
One step of the integrator for tumor positions reads
\begin{equation}
\mathbf{r}_i^{\,n+1}
   = \mathbf{r}_i^{\,n}
   + \Delta t\,[\,v_T\,\mathbf{n}_i + \mathbf{F}_i\,]
   + \sqrt{2 D_T \Delta t}\;\bm{\xi}_i^{\,n},
\label{eq:em}
\end{equation}
with $\bm{\xi}_i^n \sim \mathcal{N}(0, I_2)$ i.i.d.\ across steps and
particles. Positions are wrapped modulo $L$. Headings update by
\eqref{eq:tumor_theta}/\eqref{eq:tcell_phi} with the same Euler step.

\subsection{FTCS Euler + auto-subcycling (fields)}
The five-point Laplacian stencil on the periodic grid:
\begin{equation}
\nabla^2 c\,\big|_{i,j} =
   \frac{c_{i+1,j} + c_{i-1,j} + c_{i,j+1} + c_{i,j-1} - 4 c_{i,j}}
        {(\Delta x)^2}.
\label{eq:laplacian}
\end{equation}
The diffusion CFL ratio
$\nu = D_{\max}\Delta t / (\Delta x)^2$ must satisfy $\nu \le 1/2$ for the
explicit scheme to be stable. We use a stricter target $\nu_{\text{sub}}
\le 1/4$ and subcycle:
\begin{equation}
n_{\text{sub}} = \max\!\big(1,\,\lceil 4\nu \rceil\big),\qquad
\Delta t_{\text{sub}} = \Delta t / n_{\text{sub}}.
\end{equation}

\subsection{Gradient sampling}
$\nabla c$ is needed at off-grid T-cell positions in \eqref{eq:tcell_pos}.
We compute the central-difference gradient at every grid node and
bilinearly interpolate from the four surrounding nodes, exactly dual to
the CIC deposition~\eqref{eq:cic}.

\subsection{Minimum-image distance under PBC}
Inside every pairwise force loop:
\begin{equation}
d_x \leftarrow d_x - L\cdot\mathrm{round}(d_x/L), \qquad
d_y \leftarrow d_y - L\cdot\mathrm{round}(d_y/L).
\end{equation}

\subsection{Performance notes}
\begin{itemize}[nosep]
  \item Single \texttt{@njit(cache=True, fastmath=True)} inner step.
  \item Naive $O(N^2)$ pairwise force; at $N \lesssim 1500$ this is
        $\sim 1$~ms/step on one CPU core --- well within the time budget.
  \item Cell list ready to swap in if $N$ grows; not needed at current
        parameters.
  \item \texttt{joblib.Parallel(n\_jobs=-1, backend="loky")} for the
        192-run sweep --- $\sim$10~min wall time on 16 cores.
\end{itemize}

\section{Parameter table}

\begin{center}\small
\begin{tabular}{l l r l}
\toprule
Symbol & Code name & Value & Meaning \\
\midrule
$L$            & \texttt{L}            & 100      & box side length \\
$G$            & \texttt{G}            & 64       & field grid resolution \\
$\Delta t$     & \texttt{dt}           & 0.01     & timestep \\
$T_f$          & \texttt{T\_final}     & 100      & total simulated time \\
\midrule
$v_T$          & \texttt{v\_T}         & 0.1      & tumor self-propulsion speed \\
$D_R^T$        & \texttt{D\_R\_T}      & 0.1      & tumor rotational diffusion \\
$D_T^T$        & \texttt{D\_T\_T}      & 0.001    & tumor translational diffusion \\
$\sigma$       & \texttt{sigma\_T}     & 1.0      & repulsion cutoff (cell diameter) \\
$k_{\text{rep}}$ & \texttt{k\_rep\_T}  & 30       & repulsion stiffness \\
$p_{\text{div}}$ & \texttt{p\_div}     & 0.004    & per-step division probability \\
$r_{\text{nbr}}$ & \texttt{nbr\_radius}& 1.5      & neighbor-count radius (density gate) \\
$n_{\max}$     & \texttt{nbr\_threshold} & 6      & max neighbors allowed for division \\
\midrule
$v_I$          & \texttt{v\_I}         & 1.0      & T-cell self-propulsion speed \\
$D_R^I$        & \texttt{D\_R\_I}      & 1.0      & T-cell rotational diffusion \\
$D_T^I$        & \texttt{D\_T\_I}      & 0.001    & T-cell translational diffusion \\
$\chi_a$       & \texttt{chi\_a}       & 20.0     & attractant chemotactic coupling \\
$\alpha=\chi_s$& \texttt{chi\_s}       & scanned  & suppressant chemotactic coupling \\
\midrule
$D_a$          & \texttt{D\_a}         & 5.0      & attractant diffusion \\
$D_s$          & \texttt{D\_s}         & 0.5      & suppressant diffusion \\
$s_a$          & \texttt{s\_a}         & 1.0      & attractant secretion rate \\
$s_s$          & \texttt{s\_s}         & 1.0      & suppressant secretion rate \\
$\lambda_a$    & \texttt{lam\_a}       & 0.1      & attractant decay rate \\
$\lambda_s$    & \texttt{lam\_s}       & 0.1      & suppressant decay rate \\
\midrule
$r_{\text{kill}}$ & \texttt{r\_kill}   & 1.5      & killing engagement radius \\
$p_{\text{kill}}$ & \texttt{p\_kill}   & 0.12     & per-T-cell per-step kill probability \\
\midrule
$N_T(0)$       & \texttt{N\_T\_initial}& 50       & initial tumor seed \\
$\rho_I=N_I(0)$& \texttt{N\_I\_initial}& scanned  & initial T-cell count \\
\bottomrule
\end{tabular}
\end{center}

Five parameters were tuned away from the spec defaults so the three-phase
structure becomes visible in the 8-hour time budget: $\chi_a$ $5\to 20$,
$p_{\text{kill}}$ $0.05\to 0.12$, $p_{\text{div}}$ $0.005\to 0.004$, $T_f$
$200\to 100$, $N_T^{\max}$ $4000\to 800$. See \texttt{docs/DECISIONS.md}.

\section{Order parameter}

The primary observable is the final tumor fraction, clipped for log
plotting:
\begin{equation}
\Phi(\rho_I, \alpha) \equiv
  \operatorname{clip}\!\left(\,\frac{N_T(T_f)}{N_T(0)},\;
                              10^{-2},\; 10^{+2}\right),
\end{equation}
geometrically averaged over $S$ seeds:
\begin{equation}
\bar\Phi(\rho_I, \alpha)
  = \exp\!\Big(\,S^{-1}\sum_{s=1}^{S} \log \Phi^{(s)}(\rho_I, \alpha)\Big).
\end{equation}
Secondary observables: time-resolved $N_T(t), N_I(t)$ trajectories.
These distinguish \emph{clearance} (monotone fast decay), \emph{control /
dormancy} (long plateau at intermediate $N_T$), and \emph{escape}
(monotone rise to the array cap).

\section{Key dimensionless groups}

A way to organize the parameters:
\begin{itemize}[nosep]
  \item \textbf{Péclet of each species}:
        $\mathrm{Pe} = v / (D_R\sigma)$.\\
        $\mathrm{Pe}_T = 0.1/(0.1\cdot 1) = 1$,\quad
        $\mathrm{Pe}_I = 1.0/(1.0\cdot 1) = 1$.
  \item \textbf{Attractant range}:
        $\ell_a = \sqrt{D_a/\lambda_a} \approx 7\sigma$.
  \item \textbf{Suppressant range}:
        $\ell_s = \sqrt{D_s/\lambda_s} \approx 2.2\sigma$.
  \item \textbf{Range ratio}:
        $\ell_a/\ell_s \approx 3.2$ --- the asymmetry that enables phase
        separation.
  \item \textbf{Suppression-to-attraction ratio}:
        $\alpha/\chi_a$. At our $\chi_a = 20$, $\alpha\in[0,25]$ scans
        this ratio from 0 to 1.25 --- bracketing the regime where
        suppression overwhelms attraction.
\end{itemize}

\section{Relation to known active-matter theory}

The model lives in the same theoretical neighborhood as Toner--Tu
hydrodynamics of polar active matter (no proliferation),
Joanny--Prost--Ranft growing active matter (proliferation as a
non-conserved source), and Keller--Segel chemotaxis (the attractant
feedback loop is the classical instability). The novel content here is
the \emph{combination} of (i) long-range positive coupling via $c_a$ and
(ii) short-range negative coupling via $c_s$ from the \emph{same} tumor
source. That asymmetry generates the bistable phase boundary which maps
onto the immune-excluded tumor histology.

\section{Cell-level heterogeneity}

Replace the scalar self-propulsion $v_T$ by per-cell $v_{T,i}$ drawn from a
log-normal at birth, and similarly for rotational diffusion:
\begin{equation}
v_{T,i} \sim \mathrm{LogNormal}(\mu_v,\,\sigma_v^2),\qquad
D_{R,i}^{T} \sim \mathrm{Gamma}(k_R, \theta_R).
\end{equation}
Mean and CV are user-set. Tumor daughters inherit with a small additive
drift:
$v_{T,i'} = v_{T,i} + \varepsilon,\ \varepsilon\sim\mathcal{N}(0, \sigma_\text{daughter}^2)$.

\paragraph{Effect.} Broadens MIPS phase boundaries and creates a "slow core
+ fast invasive tail" inside the colony. Bimodal sharp phase transitions
become smoother crossovers controlled by the upper quantile of the
persistence-length distribution
$\ell_{p,i} = v_{T,i}/D_{R,i}^{T}$ (Cates \& Tailleur 2015; Maiuri et al.\ 2015).

\section{Pressure-gated proliferation (Byrne--Drasdo)}

The discrete neighbor-count gate of §2 is replaced by a continuous
mechanical-pressure gate. The local pressure is the trace of the pairwise
stress accumulated inside the existing force loop:
\begin{equation}
P_i \;=\; \frac{1}{A_i}\sum_{j\in\mathcal{N}_i} \mathbf{r}_{ij}\cdot\mathbf{F}_{ij},
\qquad A_i = \pi\sigma^2.
\end{equation}
The per-step division probability becomes
\begin{equation}
p_{\text{div},i} \;=\; p_{\text{div},0}\cdot \max\!\big(0,\,1 - P_i / P^\star\big).
\end{equation}
Daughter placement is unchanged (random direction at offset $0.3\sigma$);
no pressure-induced apoptosis is included in the current implementation.

\paragraph{Effect.} Replaces the cutoff artefact of the $n_{\max}$ rule by
a smooth homeostatic-pressure plateau. Cells at the core of a tumor stop
dividing because the local virial pressure exceeds $P^\star$; surface
cells continue. Combined with §13 below, this produces a genuinely stable
dormant fixed point: in the parameter slice where birth at the surface
balances kill, $N_T$ settles into a steady-state instead of asymptotically
escaping. References: Byrne \& Drasdo 2009; Ranft, Prost, Jülicher et al.\
2010.

\section{Macrophages: a third species with M1/M2 polarization}

Introduce a third self-propelled species $\mathbf{M}$ with positions
$\mathbf{R}_k^{(M)}$, headings $\phi_k^{(M)}$, and a continuous
polarization scalar
\begin{equation}
p_k(t) \in [-1, +1], \qquad p_k = -1: \text{M2},\ p_k = +1: \text{M1}.
\end{equation}

\subsection{Position dynamics}

Macrophages obey the same overdamped Langevin form as T cells but with
weaker chemotaxis and lower speed/persistence:
\begin{equation}
\frac{d\mathbf{R}_k^{(M)}}{dt}
  = v_M\,\mathbf{m}_k^{(M)}
  + \chi_a^{(M)}\,\nabla c_a(\mathbf{R}_k^{(M)})
  + \mathbf{F}_{\text{rep}}^{(M)}
  + \sqrt{2 D_M}\,\boldsymbol{\eta}_k^{(M)}(t),
\end{equation}
with cross-species repulsion to tumor and T cells through
$\mathbf{F}_{\text{rep}}^{(M)}$.

\subsection{Polarization dynamics}

The polarization relaxes toward an equilibrium set by local cytokine
signals:
\begin{equation}
\boxed{\;
\frac{d p_k}{dt}
  = \frac{p_{\text{eq}}(\mathbf{R}_k^{(M)}) - p_k}{\tau_p}
  + \sqrt{2 D_p}\,\zeta_k(t),
\qquad
p_{\text{eq}}(\mathbf{x}) \;=\; \tanh\!\big(b_{M_1} - \kappa_s\,c_s(\mathbf{x}) - \kappa_{IL}\,c_{IL10}(\mathbf{x})\big).
\;}
\end{equation}
$b_{M_1}$ is an external M1-promoting bias used as the treatment knob; in
the baseline run $b_{M_1}=0$ so polarization is driven entirely by local
suppressant fields (positive feedback toward M2 inside the tumor core).
$p_k$ is clipped to $[-1,+1]$ each step.

\subsection{New cytokine field: IL-10}

M2-skewed macrophages secrete a long-range CD8-repulsive cytokine,
modelled as
\begin{equation}
\frac{\partial c_{IL10}}{\partial t} \;=\; D_{IL10}\,\nabla^2 c_{IL10}
   + s_{IL10}\,\rho_{M_2}(\mathbf{x},t) - \lambda_{IL10}\,c_{IL10},
\end{equation}
\begin{equation}
\rho_{M_2}(\mathbf{x},t) \;=\; \sum_k \max(0, -p_k)\,\delta(\mathbf{x}-\mathbf{R}_k^{(M)})
   \quad \text{(only M2-skewed cells contribute)}.
\end{equation}

\subsection{Modified CD8 chemotaxis}

The CD8 equation~\eqref{eq:tcell_pos} is extended so that suppression is
no longer slaved to tumor location:
\begin{equation}
\frac{d\mathbf{R}_j}{dt}
  = v_I\,\mathbf{m}_j
  + \chi_a\,\nabla c_a(\mathbf{R}_j)
  - \alpha\,\nabla\!\left[\,c_s(\mathbf{R}_j) + c_{IL10}(\mathbf{R}_j)\,\right]
  + \mathbf{F}_{\text{rep}} + \sqrt{2 D_I}\,\boldsymbol{\zeta}_j(t).
\end{equation}

\subsection{M1 phagocytosis rule}

For each macrophage with $p_k > 0$ within $r_{\text{phag}}$ of an alive
tumor cell $i^\star$, the tumor cell is removed with probability
\begin{equation}
p_{\text{phag,eff}} \;=\; p_{\text{phag}}\cdot\frac{1 + p_k}{2}
\quad\in\;[0,\,p_{\text{phag}}].
\end{equation}
Fully M1-polarized macrophages phagocytose at rate $p_{\text{phag}}$;
fully M2 cells phagocytose at 0. At most one phagocytosis event per
macrophage per step.

\subsection{Treatment experiment: combination therapy}

The natural mid-run knobs are
\begin{itemize}[nosep]
  \item $\alpha \to \alpha'$ at $t_{\text{treat}}$ \quad (anti-PD-1 / anti-PD-L1 mimic),
  \item $b_{M_1} \to b'_{M_1} > 0$ at $t_{\text{treat}}$ \quad (TAM repolarization, e.g.\ CSF1R inhibitor or TLR agonist).
\end{itemize}
At moderate $\alpha$, anti-PD-1 alone rescues escape; TAM repolarization
alone often fails because $c_{IL10}$ has already built up and the
$-\kappa_{IL}c_{IL10}$ feedback locks polarization toward M2 even with
$b_{M_1} > 0$. Combination therapy reliably rescues escape across the
boundary region of the phase diagram.

\section{ECM density and MMP degradation}

Two additional scalar fields on the same $(G,G)$ grid:
\begin{equation}
\partial_t m \;=\; D_m\,\nabla^2 m + s_m\,\rho_T - \lambda_m\,m,
\qquad
\partial_t \rho_E \;=\; -k_{\text{deg}}\,m\,\rho_E + k_{\text{rec}}\,\rho_E(\rho_{E,0} - \rho_E).
\end{equation}
ECM density $\rho_E$ is initialized uniform at $\rho_{E,0}$ and decays
locally under tumor-secreted MMP activity; the recovery term restores
density slowly toward baseline.

The ECM couples back to tumor motility through (i) a density-dependent
drag and (ii) a pore-size cutoff:
\begin{equation}
\dot{\mathbf{r}}_i \;=\; \frac{1}{1 + \beta_{\text{drag}}\,\rho_E(\mathbf{r}_i)}\,\big[\,v_T\,\mathbf{n}_i + \mathbf{F}_i\,\big] + \sqrt{2 D_T}\,\boldsymbol{\eta}_i,
\end{equation}
with self-propulsion suppressed entirely (\,$v_{\text{eff}}=0$\,) whenever
the local pore radius
$r_p(\mathbf{x}) = r_0/\sqrt{\rho_E(\mathbf{x})}$ is smaller than the
nuclear cross-section $r_p^\star$ \emph{and} the local MMP level lies
below threshold:
\begin{equation}
v_{\text{eff}}(\mathbf{r}_i) \;=\;
  \begin{cases}
    0, & r_p(\mathbf{r}_i) < r_p^\star \text{ and } m(\mathbf{r}_i) < m^\star,\\[2pt]
    v_T/(1+\beta_{\text{drag}}\rho_E(\mathbf{r}_i)), & \text{otherwise}.
  \end{cases}
\end{equation}
The MMP gate models lytic digestion: if the local protease level exceeds
$m^\star$, the cell ``opens'' the pore and proceeds even in dense matrix.
Defaults are chosen so that $\rho_{E,0}=0$ recovers the baseline
exactly.

\paragraph{Effect.} Dense ECM compresses the escape phase and grows a
new "matrix-jammed" regime. With MMP secretion active, the tumor carves
a depleted-ECM cavity around itself --- the Friedl trail-blazing
signature. References: Wolf \& Friedl 2013; Anderson \& Chaplain 1998;
Deakin \& Chaplain 2013.

\section{Cell--cell adhesion}

A Morse-like attractive shoulder is added to the existing harmonic
repulsion, acting between tumor cells only (cadherins are an epithelial
marker):
\begin{equation}
\mathbf{F}_{cc}(\mathbf{d}) \;=\; \underbrace{k_{\text{rep}}\,\max(0,\sigma-|\mathbf{d}|)\,\hat{\mathbf{d}}}_{\text{baseline}}
   \;-\; k_{\text{adh}}\,\max(0,|\mathbf{d}|-\sigma)\cdot\max(0,\sigma_{\text{adh}}-|\mathbf{d}|)\,\hat{\mathbf{d}}.
\end{equation}
An optional Vicsek aligning torque acts on cadherin neighbors:
\begin{equation}
\frac{d\theta_i}{dt} \;\mathrel{+}\!=\; \frac{J_{\text{align}}}{|\mathcal{N}_i|}\sum_{j\in\mathcal{N}_i} \sin(\theta_j - \theta_i),
\quad \mathcal{N}_i = \{j: \sigma < |\mathbf{r}_{ij}| < \sigma_{\text{adh}}\}.
\end{equation}

\paragraph{Effect.} At $k_{\text{adh}}=0$ the tumor is a diffuse cloud
of single cells. Increasing $k_{\text{adh}}$ converts it to a cohesive
blob; adding the Vicsek term turns it into a streaming, hexagonally
packed sheet --- the morphology of collective invasion. References:
Friedl \& Gilmour 2012; Bi/Manning shape-index jamming (Bi et al.\ 2015,
Ilina et al.\ 2020).

\section{Hypoxia, angiogenesis, and extended cell repertoire}

The final extension layer adds an oxygen field, vascular ``point''
sources that move under angiogenic signalling, and three additional
immune cell populations (NK, DC, MDSC) that share the existing
overdamped-Langevin / soft-repulsion infrastructure.  Setting
\texttt{s\_O2\_vessel}, the new pool sizes, and \texttt{s\_VEGF\_hyp} to
$0$ degrades the model cleanly to the ECM-extended version of §13.

\subsection{Oxygen reaction--diffusion field}
Let $c_{O_2}(\mathbf{x},t)$ be the local oxygen concentration on the same
$(G,G)$ periodic grid.  It is sourced by discrete vessel ``nodes''
$\{\mathbf{V}_k\}_{k=1}^{N_V}$ (treated as point deposits) and consumed
by every alive cell, irrespective of species:
\begin{equation}
\partial_t c_{O_2} \;=\; D_{O_2}\,\nabla^2 c_{O_2}
   \;+\; s_{O_2}^{\text{vess}}\!\sum_{k=1}^{N_V}\delta(\mathbf{x}-\mathbf{V}_k)
   \;-\; \bigl[\lambda_{O_2} + k_{O_2}^{\text{cons}}\,\rho_{\text{all}}(\mathbf{x},t)\bigr]\,c_{O_2},
\end{equation}
with the lumped consuming density
\begin{equation}
\rho_{\text{all}}(\mathbf{x},t) \;=\; \rho_T + \rho_I + \rho_M + \rho_{NK} + \rho_{DC} + \rho_{MDSC}.
\end{equation}
The initial condition is uniform normoxia, $c_{O_2}(\mathbf{x},0) = c_{O_2,0}$.
A cell at position $\mathbf{x}$ is called \emph{hypoxic} when
$c_{O_2}(\mathbf{x}) < c_{O_2}^{\,\text{hyp}}$.  The stepper is the same
FTCS Euler subcycler used for the other RD fields, with the linear sink
term $\lambda_{O_2} + k_{O_2}^{\text{cons}}\rho_{\text{all}}$ folded into
the elementwise decay rate.

\subsection{VEGF field and angiogenic vessels}
A VEGF field is secreted only by \emph{hypoxic} tumor cells:
\begin{equation}
\partial_t c_{V} \;=\; D_V\,\nabla^2 c_V
   \;+\; s_V^{\text{hyp}}\,\rho_T^{\text{hyp}}(\mathbf{x},t)
   \;-\; \lambda_V\,c_V,
\quad
\rho_T^{\text{hyp}}(\mathbf{x},t)
   \;=\!\!\sum_{i:\,c_{O_2}(\mathbf{r}_i)<c_{O_2}^{\,\text{hyp}}}\!\!\!\!\delta(\mathbf{x}-\mathbf{r}_i).
\end{equation}
Vessels are discrete nodes that drift up the VEGF gradient with small
positional noise and may sprout to spawn new vessels:
\begin{equation}
\frac{d\mathbf{V}_k}{dt} \;=\; \chi_V\,\nabla c_V(\mathbf{V}_k)
   \;+\; \sqrt{2 D_V^{\text{pos}}}\,\bm{\eta}_k^{V}(t),
\qquad
\mathbb{P}\bigl[\text{sprout from }k \text{ in } dt\bigr]
   \;=\; r_{\text{sprout}}\,\mathbf{1}[c_V(\mathbf{V}_k)>c_V^{\star}]\,dt,
\end{equation}
with the new vessel placed at an offset of $\approx\sigma$ at a random
angle.  Vessel count is capped by $N_V^{\max}$.

\subsection{Hypoxia couplings}
Three couplings of $c_{O_2}$ back to the cell-level dynamics:
\begin{enumerate}[nosep,leftmargin=*]
  \item \textbf{O$_2$-gated proliferation.}  The pressure-gated division
        rate of §11 is multiplied by a smooth Hill-like factor in O$_2$:
        \begin{equation}
        p_{\text{div},i}^{\text{eff}} \;=\; p_{\text{div},0}\cdot\max(0,1-P_i/P^\star)\cdot
          \sigma\!\left(\frac{c_{O_2}(\mathbf{r}_i) - c_{O_2}^{\text{div}}}{c_{O_2}^{\,\text{scale}}}\right),
        \end{equation}
        with $\sigma(z)=1/(1+e^{-z})$.  Severely hypoxic tumor cells
        therefore quiesce.
  \item \textbf{VEGF secretion.}  Only hypoxic tumor cells secrete VEGF
        (see above), creating the natural feedback loop hypoxia $\to$
        VEGF $\to$ vessel drift toward the hypoxic core.
  \item \textbf{Reduced killing in hypoxic regions.}  Both CD8 and NK
        effective kill probabilities are multiplied by
        $(1 - \phi_{\text{hyp}}\,\mathbf{1}[c_{O_2}<c_{O_2}^{\,\text{hyp}}])$.
\end{enumerate}

\subsection{New cell populations}
All three new populations evolve under the standard overdamped Langevin
self-propulsion (§3) with species-specific speeds and chemotaxis
couplings.

\paragraph{Natural-killer (NK) cells.}
Innate cytotoxic effectors; faster than CD8, weak attractant chemotaxis,
\emph{not} suppressed by $c_s$ or $c_{IL10}$.  Their kill rule mirrors
the CD8 rule but with smaller $p_{\text{kill}}^{NK}$:
\begin{equation}
\frac{d\mathbf{R}_k^{NK}}{dt}
   \;=\; v_{NK}\,\mathbf{m}_k^{NK} + \chi_a^{NK}\,\nabla c_a
   + \mathbf{F}_{\text{rep}}
   + \sqrt{2 D_{NK}}\,\bm{\zeta}_k^{NK}.
\end{equation}

\paragraph{Dendritic cells (DC).}
Antigen-presenting; strong chemotaxis up $c_a$, no direct killing.
Their density is deposited onto a buff field $c_{DC}$:
\begin{equation}
\partial_t c_{DC} \;=\; D_{DC}\,\nabla^2 c_{DC} + s_{DC}\rho_{DC} - \lambda_{DC}c_{DC},
\end{equation}
which locally amplifies CD8 effective kill probability:
\begin{equation}
p_{\text{kill}}^{CD8,\text{eff}}(\mathbf{x})
   \;=\; p_{\text{kill}}\,\bigl(1 + \beta_{DC}\,c_{DC}(\mathbf{x})\bigr)\,
   \bigl(1 - \phi_{\text{hyp}}\,\mathbf{1}[c_{O_2}<c_{O_2}^{\,\text{hyp}}]\bigr).
\end{equation}

\paragraph{Myeloid-derived suppressor cells (MDSC).}
Pure suppressor myeloids: motion like macrophages (slow, low
persistence), chemotaxis up $c_a$, no killing or phagocytosis.  They
contribute an additive source to $c_s$,
\begin{equation}
\partial_t c_s \;=\; D_s\,\nabla^2 c_s
   + s_s\,\rho_T + s_s^{MDSC}\,\rho_{MDSC} - \lambda_s c_s,
\end{equation}
reinforcing the immunosuppressive halo independent of tumor location.

\subsection{Derived plotting fields}
For the renderer we compute two heuristic fields on the fly from the
simulation state (these do not feed back into the dynamics):
\begin{equation}
\text{pH}(\mathbf{x},t) \;\approx\; 7.4 - 0.9\,\tilde\rho_T(\mathbf{x}),
\qquad
\text{ROS}(\mathbf{x},t) \;\approx\; \text{ROS}_0 + A\,\tilde\rho_T(\mathbf{x})\,\bigl(1 - \tilde c_{O_2}(\mathbf{x})\bigr),
\end{equation}
with $\tilde\rho_T, \tilde c_{O_2}$ normalised to $[0,1]$.  These
reproduce the Warburg/lactate-acidified core and the
hypoxia-plus-density ROS signature familiar from TME schematics.

\section{Additional parameters}

Defaults for heterogeneity, pressure-gated division, macrophages,
ECM/MMP, cadherin adhesion, and the TME-extension layer.  Setting any
of $\rho_{E,0}$, $s_m$, $k_{\text{adh}}$, $J_{\text{align}}$, the
macrophage seed, or the new NK/DC/MDSC seed counts to $0$ disables the
corresponding mechanism cleanly.

\begin{center}\small
\begin{tabular}{l l r l}
\toprule
Symbol & Code name & Value & Meaning \\
\midrule
\multicolumn{4}{l}{\emph{Pressure-gated proliferation}}\\
$P^\star$           & \texttt{P\_star}      & 8.0   & homeostatic pressure scale \\
$p_{\text{div},0}$  & \texttt{p\_div0}      & 0.004 & nominal per-step division probability \\
\midrule
\multicolumn{4}{l}{\emph{Macrophages \& IL-10 field}}\\
$v_M$               & \texttt{v\_M\_mean}   & 0.2   & macrophage self-propulsion speed \\
$\sigma_M$          & \texttt{sigma\_M}     & 1.0   & macrophage repulsion cutoff \\
$k_{\text{rep}}^{M}$& \texttt{k\_rep\_M}    & 30    & macrophage repulsion stiffness \\
$\chi_a^{(M)}$      & \texttt{chi\_a\_M}    & 8.0   & macrophage attractant coupling \\
$\tau_p$            & \texttt{tau\_p}       & 15.0  & polarization relaxation time \\
$D_p$               & \texttt{D\_p}         & 0.03  & polarization noise \\
$\kappa_s$          & \texttt{kappa\_s}     & 8.0   & suppressant $\to$ M2 drive \\
$\kappa_{IL}$       & \texttt{kappa\_il}    & 4.0   & IL-10 $\to$ M2 positive feedback \\
$b_{M_1}$           & \texttt{M1\_bias}     & 0.0   & external M1-bias (treatment knob) \\
$r_{\text{phag}}$   & \texttt{r\_phag}      & 1.8   & M1 phagocytosis engagement radius \\
$p_{\text{phag}}$   & \texttt{p\_phag}      & 0.10  & base phagocytosis probability \\
$D_{IL10}$          & \texttt{D\_IL10}      & 4.0   & IL-10 diffusion (long-range) \\
$s_{IL10}$          & \texttt{s\_IL10}      & 1.0   & IL-10 secretion rate per M2 \\
$\lambda_{IL10}$    & \texttt{lam\_IL10}    & 0.1   & IL-10 decay rate \\
\midrule
\multicolumn{4}{l}{\emph{ECM density + MMP}}\\
$\rho_{E,0}$        & \texttt{rho\_E\_init} & 0 (off) & baseline ECM density (set $>0$ to enable) \\
$D_m$               & \texttt{D\_m}         & 3.0   & MMP diffusion \\
$s_m$               & \texttt{s\_m}         & 0 (off) & tumor MMP secretion rate \\
$\lambda_m$         & \texttt{lam\_m}       & 0.15  & MMP decay rate \\
$k_{\text{deg}}$    & \texttt{k\_deg}       & 1.0   & MMP-driven ECM degradation rate \\
$k_{\text{rec}}$    & \texttt{k\_rep\_ECM}  & 0.005 & logistic ECM recovery rate \\
$\beta_{\text{drag}}$ & \texttt{beta\_drag} & 0.8   & ECM drag coefficient \\
$r_0$               & \texttt{r\_0}         & 1.0   & baseline pore-radius constant \\
$r_p^\star$         & \texttt{r\_p\_star}   & 0.55  & nuclear-cross-section pore threshold \\
$m^\star$           & \texttt{mmp\_open\_thresh} & 0.10 & MMP level opening jammed pores \\
\midrule
\multicolumn{4}{l}{\emph{Cell--cell adhesion (cadherins)}}\\
$k_{\text{adh}}$    & \texttt{k\_adh}       & 8.0   & attractive shoulder stiffness \\
$\sigma_{\text{adh}}$ & \texttt{sigma\_adh} & 1.6   & adhesion outer cutoff \\
$J_{\text{align}}$  & \texttt{J\_align}     & 0.5   & Vicsek alignment coupling \\
\midrule
\multicolumn{4}{l}{\emph{Oxygen + VEGF + vasculature}}\\
$D_{O_2}$              & \texttt{D\_O2}              & 8.0   & oxygen diffusion \\
$\lambda_{O_2}$        & \texttt{lam\_O2}            & 0.02  & oxygen natural decay \\
$k_{O_2}^{\text{cons}}$& \texttt{k\_O2\_cons}        & 0.35  & per-cell oxygen consumption \\
$c_{O_2,0}$            & \texttt{O2\_init}           & 0.85  & uniform initial O$_2$ \\
$c_{O_2}^{\,\text{hyp}}$ & \texttt{O2\_hyp\_thresh}  & 0.45  & hypoxia threshold \\
$c_{O_2}^{\text{div}}$ & \texttt{O2\_div\_thresh}    & 0.5   & O$_2$ midpoint of division gate \\
$c_{O_2}^{\,\text{scale}}$ & \texttt{O2\_div\_scale} & 0.10  & sigmoid sharpness \\
$\phi_{\text{hyp}}$    & \texttt{hypoxia\_kill\_penalty} & 0.6 & hypoxic-region kill penalty \\
$s_{O_2}^{\text{vess}}$& \texttt{s\_O2\_vessel}      & 3.0   & per-vessel O$_2$ deposit rate \\
$\chi_V$               & \texttt{chi\_vessel}        & 80.0  & vessel drift up VEGF grad \\
$D_V^{\text{pos}}$     & \texttt{D\_vessel}          & 0.02  & vessel positional noise \\
$N_V^{\,0}$            & \texttt{n\_vessels\_init}   & 6     & initial vessel count \\
$N_V^{\max}$           & \texttt{n\_vessels\_max}    & 32    & maximum vessel count (sprouting cap) \\
$r_{\text{sprout}}$    & \texttt{sprout\_rate}       & 0.03  & per-step sprout probability \\
$c_V^{\star}$          & \texttt{sprout\_VEGF\_thresh} & 0.005 & VEGF threshold to allow sprout \\
$D_V$                  & \texttt{D\_VEGF}            & 3.0   & VEGF diffusion \\
$\lambda_V$            & \texttt{lam\_VEGF}          & 0.10  & VEGF decay \\
$s_V^{\text{hyp}}$     & \texttt{s\_VEGF\_hyp}       & 1.0   & VEGF source per hypoxic tumor cell \\
\midrule
\multicolumn{4}{l}{\emph{NK cells (innate cytotoxic)}}\\
$v_{NK}$               & \texttt{v\_NK\_mean}        & 1.2   & NK self-propulsion speed \\
$D_R^{NK}$             & \texttt{D\_R\_NK\_mean}     & 0.8   & NK rotational diffusion \\
$\chi_a^{NK}$          & \texttt{chi\_a\_NK}         & 8.0   & NK attractant chemotaxis \\
$r_{\text{kill}}^{NK}$ & \texttt{r\_kill\_NK}        & 1.4   & NK kill engagement radius \\
$p_{\text{kill}}^{NK}$ & \texttt{p\_kill\_NK}        & 0.07  & NK per-step kill probability \\
\midrule
\multicolumn{4}{l}{\emph{Dendritic cells (DC) + buff field $c_{DC}$}}\\
$v_{DC}$               & \texttt{v\_DC\_mean}        & 0.6   & DC self-propulsion speed \\
$\chi_a^{DC}$          & \texttt{chi\_a\_DC}         & 30.0  & DC attractant chemotaxis (strong) \\
$D_{DC}$               & \texttt{D\_DC}              & 1.0   & DC-buff field diffusion (short range) \\
$s_{DC}$               & \texttt{s\_DC}              & 1.0   & DC-buff field source per DC \\
$\lambda_{DC}$         & \texttt{lam\_DC}            & 0.3   & DC-buff field decay \\
$\beta_{DC}$           & \texttt{cd8\_buff\_scale}   & 4.0   & DC-buff multiplier on $p_{\text{kill}}^{CD8}$ \\
\midrule
\multicolumn{4}{l}{\emph{MDSC (suppressor myeloids)}}\\
$v_{MDSC}$             & \texttt{v\_MDSC\_mean}      & 0.2   & MDSC self-propulsion speed \\
$\chi_a^{MDSC}$        & \texttt{chi\_a\_MDSC}       & 6.0   & MDSC attractant chemotaxis \\
$s_s^{MDSC}$           & \texttt{s\_s\_MDSC}         & 0.6   & MDSC contribution to $c_s$ source \\
\bottomrule
\end{tabular}
\end{center}

\section{Roadmap of further mechanisms}

The following are flagged in \texttt{docs/extensions.md} with equations
ready, not yet implemented:

\begin{itemize}[nosep]
  \item \textbf{§2 cell--ECM adhesion (integrins).} Per-cell engaged-integrin
        state $e_i$ with biphasic DiMilla speed law.
  \item \textbf{§4 ECM stiffness / durotaxis.} Young's-modulus field $E_m$
        with drift $\chi_E\nabla E_m$ and stiffness-dependent speed.
  \item \textbf{§5 fiber alignment.} Nematic director $Q(\mathbf{x})$ with
        contact-guidance torque; anisotropic mobility tensor.
  \item \textbf{Tregs, NK cells, MDSCs, CAFs.} See \texttt{docs/extensions.md}
        Part~C for the active-matter representation of each.
  \item \textbf{EMT continuous state.} Bistable Landau-like potential
        coupling motility, adhesion, and MMP secretion to a single
        epithelial-mesenchymal coordinate $s_i \in [0,1]$.
\end{itemize}

\section{Headline result}

\begin{quote}
\emph{The transition between immune clearance and tumor escape is sharp
and bistable in this active matter. A simulated checkpoint inhibitor ---
a sudden reduction of $\alpha$ at fixed $\rho_I$ --- moves the same
initial state from escape to clearance, demonstrating that checkpoint
inhibition is mechanistically equivalent to crossing the phase boundary
in this model. Macrophages make the suppressant field move and adapt:
$c_{IL10}$ is sourced by a non-tumor species and follows the immune
infiltrate, breaking the spatial slaving of suppression to tumor
density, so that combination therapy --- simultaneous $\alpha$ knockdown
and $M_1$-bias activation --- crosses the phase boundary even when
either monotherapy fails. ECM density together with MMP secretion
produces a matrix-jammed regime in which the tumor digests a depleted
cavity around itself (Friedl trail-blazing). Cadherin adhesion plus
Vicsek alignment converts a diffuse motile tumor into a streaming
epithelial sheet. Pressure-gated proliferation produces a genuinely
stable dormant fixed point in the boundary region of the phase diagram,
populating the ``control'' phase that was only transient in the
minimal two-species model.}
\end{quote}

\end{document}
